\dykband{EXERCISES IN SWAFL}
{\small Open SWAFL (desktop or online) with the benchmark project. The chapter of the \href{https://github.com/halilunaziru73-creator/Coupled-Water-Fluxes-Pore-Scale-Dynamics-Precision-Irrigation-Mulching-Feedback-Loops/raw/main/Software/v1.1.0/SWAFL_Desktop/docs/manual_v1.1.0/SWAFL_Manual_v1.1.0.pdf}{SWAFL manual} that explains each task is given in brackets; record your values in a table and compare them with what is expected (\emph{Check}).}\par\vspace{3pt}
\begin{exlist}
\item {\color{khblue}[Manual, Chapter 8]} Open \emph{Tillage and structure}; compare no-till and tillage for the loam and record the drainage and the evaporation.\par\noindent{\small\itshape Check: tillage raises the infiltration and changes the drainage.}
\item {\color{khblue}[Manual, Chapter 9]} Open \emph{Mulch control}; vary the mulch strength from 0 to the maximum and record the saving of evaporation.\par\noindent{\small\itshape Check: the saving rises and then levels off.}
\item {\color{khblue}[Manual, Chapter 9]} Open \emph{Fate of saved water} and \emph{Regime map}; find the rain scale at which most of the saving drains.\par\noindent{\small\itshape Check: in wet climates most of the saving drains.}
\item {\color{khblue}[Manual, Chapter 10]} Open \emph{Feedback loops}; follow the seven effects and write the sentence under the picture in your own words.\par\noindent{\small\itshape Check: each effect has a sign and a size.}
\end{exlist}
